\documentclass[12pt, a4paper]{article}
\usepackage{../../../myconfig} % Gọi file cấu hình từ ngoài gốc vào

% ==========================================
% BẮT ĐẦU NỘI DUNG TÀI LIỆU
% ==========================================
\begin{document}

% Truyền 3 tham số: {Nhãn cam} {Chữ to chính giữa} {Chữ ở góc phải Header}
\titlebox{Chủ đề: Hàm số}{Cấp số cộng- cấp số nhân}{Hàm số: CSC-CSN}

% --- CÂU 1 ---
\questionLinesManual{0}{1}{Cho cấp số cộng \( (u_n) \) với \( u_1 = 2 \) và công sai \( d = 3 \). Giá trị của \( u_5 \) bằng}{
    \begin{tasks}(4)
        \task $14$
        \task $17$
        \task $15$
        \task $12$
    \end{tasks}
}

% --- CÂU 2 ---
\questionLinesManual{0}{2}{Cấp số cộng \( (u_n) \) có \( u_1 = 1 \) và \( u_2 = 3 \). Số hạng \( u_5 \) của cấp số cộng là:}{
    \begin{tasks}(4)
        \task $5$
        \task $7$
        \task $9$
        \task $11$
    \end{tasks}
}

% --- CÂU 3 ---
\questionLinesManual{0}{3}{Cho cấp số cộng \( (u_n) \) có \( u_1 = -2 \) và công sai \( d = 3 \). Tìm số hạng \( u_{10} \).}{
    \begin{tasks}(4)
        \task \( u_{10} = -2 \cdot 3^9 \)
        \task \( u_{10} = 25 \)
        \task \( u_{10} = -29 \)
        \task \( u_{10} = 28 \)
    \end{tasks}
}

% --- CÂU 4 ---
\questionLinesManual{0}{4}{Cho cấp số nhân \( (u_n) \) với \( u_1 = 6 \) và \( u_2 = -12 \). Công bội \( q \) của cấp số nhân đã cho là}{
    \begin{tasks}(4)
        \task \( q = -\dfrac{1}{2} \)
        \task \( q = -2 \)
        \task \( q = -18 \)
        \task \( q = -6 \)
    \end{tasks}
}

% --- CÂU 5 ---
\questionLinesManual{0}{5}{Cho cấp số cộng có số hạng đầu \( u_1 = 3 \), công sai \( d = -4 \). Số hạng thứ năm của cấp số cộng là.}{
    \begin{tasks}(4)
        \task $768$
        \task $-13$
        \task $-3072$
        \task $-17$
    \end{tasks}
}
% --- CÂU 6 ---
\questionLinesManual{0}{6}{Cho cấp số cộng \( (u_n) \), biết \( u_1 = 1 \), \( d = 2 \). Giá trị của \( u_{15} \) bằng:}{
    \begin{tasks}(4)
        \task $31$
        \task $29$
        \task $35$
        \task $27$
    \end{tasks}
}


% --- CÂU 7 ---
\questionLinesManual{0}{7}{Cho cấp số cộng \( (u_n) \) có công sai \( d = 2 \) và tổng của 25 số hạng đầu là 625. Tính \( u_5 \).}{
    \begin{tasks}(4)
        \task \( u_5 = 16 \)
        \task \( u_5 = 9 \)
        \task \( u_5 = 12 \)
        \task \( u_5 = 22 \)
    \end{tasks}
}

% --- CÂU 8 ---
\questionLinesManual{0}{8}{Cho cấp số cộng \( (u_n) \) biết \( u_1 = 2 \), công sai \( d = -5 \). Tổng 10 số hạng đầu của cấp số cộng đó là}{
    \begin{tasks}(4)
        \task $-410$
        \task $-205$
        \task $245$
        \task $-230$
    \end{tasks}
}

% --- CÂU 9 ---
\questionLinesManual{0}{9}{Cho cấp số nhân \( (u_n) \) có số hạng đầu \( u_1 = 2 \) và \( u_6 = 486 \). Công bội \( q \) bằng}{
    \begin{tasks}(4)
        \task \( q = 5 \)
        \task \( q = \dfrac{3}{2} \)
        \task \( q = \dfrac{2}{3} \)
        \task \( q = 3 \)
    \end{tasks}
}

% --- CÂU 10 ---
\questionLinesManual{0}{10}{Cho cấp số cộng \( (u_n) \) có \( u_4 = -12 \), \( u_{14} = 18 \). Tính tổng 18 số hạng đầu tiên của cấp số cộng này.}{
    \begin{tasks}(4)
        \task \( S_{18} = -72 \)
        \task \( S_{18} = 96 \)
        \task \( S_{18} = -96 \)
        \task \( S_{18} = 72 \)
    \end{tasks}
}

% --- CÂU 11 ---
\questionLinesManual{0}{11}{Cho cấp số cộng \( (u_n) \) có \( u_1 = -2 \) và công sai \( d = 3 \). Tìm số hạng tổng quát \( u_n \) của cấp số cộng đó.}{
    \begin{tasks}(4)
        \task \( u_n = 3n + 5 \)
        \task \( u_n = 3n - 5 \)
        \task \( u_n = 3n - 1 \)
        \task \( u_n = 3n + 1 \)
    \end{tasks}
}

% --- CÂU 12 ---
\questionLinesManual{0}{12}{Dãy số nào sau đây không phải là cấp số nhân?}{
    \begin{tasks}(4)
        \task $1; 2; 3; 4; 5$
        \task $1; 2; 4; 8; 16$
        \task $1; -1; 1; -1; 1$
        \task $1; -2; 4; -8; 16$
    \end{tasks}
}

% --- CÂU 13 ---
\questionLinesManual{0}{13}{Cho cấp số nhân \( (u_n) \) có \( u_1 = 3 \) và \( q = 2 \). Khi đó, số hạng tổng quát của cấp số nhân là}{
    \begin{tasks}(4)
        \task \( u_n = 2 \cdot 3^{n-1} \)
        \task \( u_n = 3 \cdot 2^{n-1} \)
        \task \( u_n = 2 \cdot 3^n \)
        \task \( u_n = 3 \cdot 2^n \)
    \end{tasks}
}

% --- CÂU 14 ---
\questionLinesManual{0}{14}{Cho \( (u_n) \) là cấp số nhân, công bội \( q > 0 \). Biết \( u_1 = 1, u_3 = 4 \). Tìm \( u_4 \).}{
    \begin{tasks}(4)
        \task \( \dfrac{11}{2} \)
        \task $2$
        \task $16$
        \task $8$
    \end{tasks}
}

% --- CÂU 15 ---
\questionLinesManual{0}{15}{Cho một cấp số cộng \( (u_n) \) có \( u_1 = 5 \) và \( S_{50} = 5150 \). Tìm \( S_{20} \)?}{
    \begin{tasks}(4)
        \task $860$
        \task $900$
        \task $1080$
        \task $1030$
    \end{tasks}
}

% --- CÂU 16 ---
\questionLinesManual{0}{16}{Cho cấp số nhân \( (u_n) \), biết \( u_1 = 2; q = 3 \) có số hạng tổng quát là:}{
    \begin{tasks}(4)
        \task \( u_n = 2 \cdot 2^{n-1} \)
        \task \( u_n = 2 \cdot 3^{n-1} \)
        \task \( u_n = 2 \cdot 3^{n-2} \)
        \task \( u_n = 3 \cdot 2^{n-1} \)
    \end{tasks}
}

% --- CÂU 17 ---
\questionLinesManual{0}{17}{Cho cấp số nhân \( (u_n) \) với \( u_1 = 3, u_4 = 24 \). Tìm số hạng thứ 5 của cấp số nhân.}{
    \begin{tasks}(4)
        \task \( S_5 = 93 \)
        \task \( S_5 = -93 \)
        \task \( u_5 = 48 \)
        \task \( u_5 = 27 \)
    \end{tasks}
}

% --- CÂU 18 ---
\questionLinesManual{0}{18}{Cho cấp số nhân \( (u_n) \) có số hạng đầu \( u_1 = 5 \) và công bội \( q = -2 \). Số hạng thứ sáu của \( (u_n) \) là:}{
    \begin{tasks}(4)
        \task \( u_6 = -320 \)
        \task \( u_6 = 320 \)
        \task \( u_6 = 160 \)
        \task \( u_6 = -160 \)
    \end{tasks}
}

% --- CÂU 19 ---
\questionLinesManual{0}{19}{Cho cấp số nhân \( (u_n) \) biết \( u_n = 2^n, \forall n \in \mathbb{N}^* \). Tìm số hạng đầu \( u_1 \) và công bội \( q \) của cấp số nhân trên:}{
    \begin{tasks}(4)
        \task \( u_1 = 2; \, q = -2 \)
        \task \( u_1 = -2; \, q = 2 \)
        \task \( u_1 = 1; \, q = 2 \)
        \task \( u_1 = 2; \, q = 2 \)
    \end{tasks}
}

% --- CÂU 20 ---
\questionLinesManual{0}{20}{Cho cấp số nhân \( (u_n) \) có số hạng đầu \( u_1 = 5 \) và công bội \( q = -2 \). Tính \( S_6 \).}{
    \begin{tasks}(4)
        \task \( S_6 = -\dfrac{155}{3} \)
        \task \( S_6 = -\dfrac{315}{3} \)
        \task \( S_6 = -315 \)
        \task \( S_6 = 315 \)
    \end{tasks}
}

% --- CÂU 21 ---
\questionLinesManual{0}{21}{Một cấp số nhân có số hạng thứ hai bằng 4 và số hạng thứ sáu bằng 64, thì số hạng tổng quát của cấp số nhân đó có thể tính theo công thức nào dưới đây}{
    \begin{tasks}(4)
        \task \( u_n = 2^{n-1} \)
        \task \( u_n = 2^n \)
        \task \( u_n = 2^{n+1} \)
        \task \( u_n = 2n \)
    \end{tasks}
}

% --- CÂU 22 ---
\questionLinesManual{0}{22}{Cho cấp số nhân \( (u_n) \) có \( u_1 = 3 \) và \( q = -4 \). Số 768 là số hạng thứ bao nhiêu của cấp số nhân đã cho?}{
    \begin{tasks}(4)
        \task Số hạng thứ 5
        \task Số hạng thứ 6
        \task Số hạng thứ 7
        \task Không là số hạng của cấp số đã cho
    \end{tasks}
}

% --- CÂU 23 ---
\questionLinesManual{0}{23}{Tìm \( x \) để các số 2; 8; \( x \) theo thứ tự đó lập thành một cấp số nhân.}{
    \begin{tasks}(4)
        \task \( x = 14 \)
        \task \( x = 32 \)
        \task \( x = 64 \)
        \task \( x = 68 \)
    \end{tasks}
}

% --- CÂU 24 ---
\questionLinesManual{0}{24}{Cho một cấp số cộng \( (u_n) \) có số hạng tổng quát là \( u_n = 3n - 5 \). Công sai \( d \) của cấp số cộng trên là:}{
    \begin{tasks}(4)
        \task \( d = 3 \)
        \task \( d = 5 \)
        \task \( d = -5 \)
        \task \( d = -3 \)
    \end{tasks}
}

% --- CÂU 25 ---
\questionLinesManual{0}{25}{Cho một cấp số cộng \( (u_n) \) có số hạng thứ tự bằng \(-12\) và số hạng thứ mười bốn bằng 18. Tính tổng 16 số hạng đầu tiên của cấp số cộng \( (u_n) \).}{
    \begin{tasks}(4)
        \task $-24$
        \task $25$
        \task $26$
        \task $24$
    \end{tasks}
}

% --- CÂU 26 ---
\questionLinesManual{0}{26}{Cho một cấp số nhân \( (u_n) \) có \( u_1 = 2 \) và \( q = -3 \). Số 162 là số hạng thứ mấy của cấp số nhân đã cho?}{
    \begin{tasks}(4)
        \task Số hạng thứ 2
        \task Số hạng thứ 3
        \task Số hạng thứ 4
        \task Số hạng thứ 5
    \end{tasks}
}

% --- CÂU 27 ---
\questionLinesManual{0}{27}{Cho một cấp số cộng có 15 số hạng, biết tổng các số hạng của cấp số cộng bằng 225, số hạng cuối bằng 29. Tìm số hạng đầu \( u_1 \) của cấp số cộng đó.}{
    \begin{tasks}(4)
        \task \( u_1 = 5 \)
        \task \( u_1 = 3 \)
        \task \( u_1 = 1 \)
        \task \( u_1 = 2 \)
    \end{tasks}
}

% --- CÂU 28 ---
\questionLinesManual{0}{28}{Cho cấp số cộng có \( u_1 = 16 \), \( u_8 = 2 \). Tính \( S_8 \).}{
    \begin{tasks}(4)
        \task \( S_8 = 72 \)
        \task \( S_8 = 144 \)
        \task \( S_8 = 56 \)
        \task \( S_8 = 18 \)
    \end{tasks}
}

% --- CÂU 29 ---
\questionLinesManual{0}{29}{Trong các dãy số sau, dãy số nào là cấp số nhân?}{
    \begin{tasks}(4)
        \task $-1; \dfrac{1}{3}; \dfrac{1}{9}; -9; \ldots$
        \task $2; 4; 6; 8; 9; \ldots$
        \task $1; \dfrac{1}{3}; \dfrac{1}{9}; \dfrac{1}{27}; \dfrac{1}{81}; \ldots$
        \task $5; 10; 15; 20; 25; \ldots$
    \end{tasks}
}

% --- CÂU 30 ---
\questionLinesManual{0}{30}{Cho dãy số \( (u_n) \) với \( u_n = 3^n \). Số hạng \( u_{n+1} \) bằng}{
    \begin{tasks}(4)
        \task \( 3^n + 1 \)
        \task \( 3^n + 3 \)
        \task \( 3^n \cdot 3 \)
        \task \( 3(n + 1) \)
    \end{tasks}
}

% --- CÂU 31 ---
\questionLinesManual{0}{31}{Các số \(-3, -9, a\) theo thứ tự đó lập thành một cấp số nhân. Khi đó \( a \) bằng bao nhiêu?}{
    \begin{tasks}(4)
        \task \( a = 27 \)
        \task \( a = -27 \)
        \task \( a = -3 \)
        \task \( a = 3 \)
    \end{tasks}
}

% --- CÂU 32 ---
\questionLinesManual{0}{32}{Xác định công sai của cấp số cộng sau: \( 2; 5; 8; 11; 14; \ldots \)}{
    \begin{tasks}(4)
        \task $2$
        \task $3$
        \task $5$
        \task $8$
    \end{tasks}
}

% --- CÂU 33 ---
\questionLinesManual{0}{33}{Cho cấp số nhân, có ba số hạng liên tiếp là \( x - 2 \), \( x \), \( x + 4 \). Công bội của cấp số nhân đó là?}{
    \begin{tasks}(4)
        \task $1$
        \task $2$
        \task $-1$
        \task $-2$
    \end{tasks}
}

% --- CÂU 34 ---
\questionLinesManual{0}{34}{Trong các dãy số sau đây, dãy số nào là cấp số cộng?}{
    \begin{tasks}(4)
        \task \( u_n = 3n^2 + 2025 \)
        \task \( u_n = 3n + 2024 \)
        \task \( u_n = 3^n \)
        \task \( u_n = (-3)^{n+1} \)
    \end{tasks}
}

% --- CÂU 35 ---
\questionLinesManual{0}{35}{Cho cấp số cộng \( (u_n) \) có số hạng tổng quát là \( u_n = 3n - 2 \). Tìm công sai \( d \) của cấp số cộng.}{
    \begin{tasks}(4)
        \task \( d = 3 \)
        \task \( d = 2 \)
        \task \( d = -2 \)
        \task \( d = -3 \)
    \end{tasks}
}

% --- CÂU 36 ---
\questionLinesManual{0}{36}{Trong các dãy số sau đây, dãy số nào là một cấp số cộng?}{
    \begin{tasks}(2)
        \task \( u_n = n^2 + 1, n \geq 1 \)
        \task \( u_n = 2^n, n \geq 1 \)
        \task \( u_n = \sqrt{n+1}, n \geq 1 \)
        \task \( u_n = 2n - 3, n \geq 1 \)
    \end{tasks}
}

% --- CÂU 37 ---
\questionLinesManual{0}{37}{Trong các dãy số \( (u_n) \) sau đây, dãy số nào là cấp số nhân?}{
    \begin{tasks}(4)
        \task \( u_n = 3n \)
        \task \( u_n = 2^n \)
        \task \( u_n = \dfrac{1}{n} \)
        \task \( u_n = 2^n + 1 \)
    \end{tasks}
}

% --- CÂU 38 ---
\questionLinesManual{0}{38}{Cho cấp số nhân có các số hạng lần lượt là \( \dfrac{1}{4} \); \( \dfrac{1}{2} \); 1; ...; 2048. Tính tổng \( S \) của tất cả các số hạng của cấp số nhân đã cho.}{
    \begin{tasks}(4)
        \task \( S = 2047,75 \)
        \task \( S = 2049,75 \)
        \task \( S = 4095,75 \)
        \task \( S = 4096,75 \)
    \end{tasks}
}

% --- CÂU 39 ---
\questionLinesManual{0}{39}{Cho cấp số cộng hữu hạn \( (u_n) : 3, \, a, \, 13, \, b \). Tích \( ab \) bằng}{
    \begin{tasks}(4)
        \task $144$
        \task $39$
        \task $26$
        \task $104$
    \end{tasks}
}

% --- CÂU 40 ---
\questionLinesManual{0}{40}{Biết \( x \) thỏa mãn \( x^2 - 2, \, x, \, 5 - 6x \) lập thành cấp số cộng. Tính tổng bình phương các giá trị \( x \) tìm được.}{
    \begin{tasks}(4)
        \task $12$
        \task $17$
        \task $26$
        \task $10$
    \end{tasks}
}


\end{document}